\section{工程调试与问题复盘}
\label{sec:debug}

本项目的开发过程里，最有价值的部分不只是功能逐步增加，也包括一些
真正改变设计的调试经历。本章保留几个对系统形态影响较大的问题，重点说明
它们如何影响后续取舍。

\subsection{语音链路的底层排查}

语音是项目中最早让系统“像教练”的模块，也是最容易误判根因的模块。起初
使用摄像头自带的数字音频输入，延迟和识别率都不理想；后来换到板端麦克风
链路后，又遇到唤醒词不稳定、录音为空、ALSA 设备被占用等问题。

这些问题一开始很容易被归因到语音识别代码。但单独跑录音测试后可以发现，
真正麻烦的是音频设备句柄、采样率初始化、旧守护进程残留和 TTS/录音抢
设备等问题交织在一起。于是项目后期把语音调试拆成更小的链路：先确认录音
文件，再确认 ASR 返回，再进入唤醒词、命令路由和 TTS 播报。

这次经历带来的教训是，嵌入式项目里不能只在最终业务模块里查问题。音频、
摄像头、网络和传感器都需要先有最小可复现测试。只有确认底层采集链路是
干净的，再回到业务逻辑里看状态机和路由，才不会反复怀疑错误对象。

\subsection{肌电输入的佩戴影响}

肌电问题也有类似特点。早期波形不稳定时，直觉上会怀疑电源、放大电路或
代码解析。但实际测试中，贴片是否贴紧、皮肤是否潮湿、线缆是否晃动，都会
影响采样结果。贴片接触不好时，模型端再怎么调整也很难得到稳定判断。

这也是项目后来强调规范穿戴流程和 MVC 归一化的原因。模型再轻量、代码再
整洁，如果输入本身漂移很大，最终判断仍然会不稳定。与其在模型端不断补
规则，不如先把佩戴、采样和归一化流程固定下来。

\subsection{功能叠加后的重新分层}

项目后期另一个明显问题是代码耦合。每加一个功能，最简单的做法都是在旧的
主循环或前端逻辑里再加一段 if/else。但这种做法积累到一定程度后，调试
成本会快速上升：一个语音命令可能同时影响前端、FSM、LLM、TTS 和状态文件，
任何一处插入不当都会引发连锁问题。

V7.30 前后的整理没有推翻整个系统，而是把几个高频冲突点拆开：语音状态、
turn 协议、DeepSeek client、tools/router、auto trigger 和后台记录接口。
整理后代码不一定立刻变少，但入口更清楚，后续手测、拍摄和复盘也更容易
定位问题。

\subsection{调试工作台的价值}

随着模块增多，单靠终端日志很难支撑排查。第 \ref{sec:implementation} 章
描述的端到端调试控制台在工程意义上的价值，是把"系统当前处于哪一个
状态"从一个需要查日志和源码才能回答的问题，变成一个查面板就能回答
的问题。

这种工作台思路也改变了报告和演示的组织方式。过去更容易只讲“功能能做
什么”，现在可以讲清楚“系统如何被维护”。对于一个包含视觉、肌电、语音、
LLM 和 Web 的嵌入式项目来说，可维护性本身就是完成度的一部分。

